\section{Introduction}

With the continuous expansion of network scale and complexity, traditional network architectures are increasingly challenged in terms of flexibility, scalability, and automated management. 
Software-Defined Networking (SDN)~\cite{JANG2023105736} has emerged as a novel architectural paradigm that decouples the control plane from the data plane, enabling centralized management and flexible configuration of network resources~\cite{9360650}. 
While SDN's centralized control plane enables dynamic network orchestration through protocols such as OpenFlow~\cite{10.1145/1355734.1355746}, effective management depends not merely on rule deployment but also on real-time state awareness. 
% Accurate and timely acquisition of network states is thus indispensable, underscoring monitoring as a fundamental SDN capability~\cite{10.1145/3556973}.
\rev{Accurate and timely acquisition of network states is essential, making monitoring a fundamental SDN capability~\cite{10.1145/3556973}.}

SDN switches inherently support the collection of flow-level statistics such as packet counts, byte volumes, and flow duration, providing a foundation for fine-grained traffic monitoring \cite{KALAMBE2025104081}. Existing monitoring mechanisms primarily fall into two categories: push-based \cite{10.1145/3373360.3380830, 6888899, 10.1007/978-3-642-36516-4_4} and pull-based \cite{7037094, CASTILLO2020107108, 8411125, 10.1007/978-3-642-12334-4_21} approaches. 
Push-based schemes periodically or event-drivenly transmit flow statistics to the controller, offering better timeliness but often leading to congestion in the control channel and excessive load on the controller. 
In contrast, pull-based schemes rely on the controller to actively request statistics, offering greater flexibility but suffering from trade-offs between accuracy and overhead: short polling intervals increase system burden, while longer intervals risk delayed anomaly detection. 
Moreover, in either case, the centralized accumulation of monitoring data at the controller may introduce performance bottlenecks, undermining critical tasks such as route computation and flow table updates \cite{8291608}.

The advent of data plane programmability, particularly with the rise of the P4 (Programming Protocol-Independent Packet Processors) language~\cite{10.1145/2656877.2656890}, has revolutionized traffic monitoring in SDN. 
P4 enables protocol-agnostic and behavior-customized packet processing, allowing switches to implement more flexible and efficient data collection strategies \cite{HAUSER2023103561}. 
For instance, by embedding lightweight filtering or labeling logic within the switch pipeline, non-essential data can be preprocessed and discarded locally, thereby reducing the volume of traffic forwarded to the controller and alleviating pressure on the control plane \cite{10130445}. 
Nonetheless, designing a monitoring mechanism that achieves high accuracy with low overhead under a programmable data plane (PDP)~\cite{PDP} architecture remains an open and challenging problem \cite{KAUR2021109}.

Meanwhile, Machine Learning (ML) has shown considerable promise in network management, particularly in behavior modeling of complex systems, anomaly detection, and policy optimization \cite{JAMSHIDI2025111103, JISI2025110939, Serag2025, 10480735, 10273389, 10198239, FOSIC2023100466, 9505258}. 
In the context of SDN monitoring, ML-based approaches can assist the controller in dynamically adjusting monitoring policies under uncertain or time-varying traffic conditions---e.g., by selectively sampling flows \cite{9335605}---thus improving the balance between monitoring accuracy and system overhead. 
Although prior works have investigated Reinforcement Learning (RL)~\cite{RL}, Supervised Learning (SL)~\cite{SupervisedLearning}, or Federated Learning (FL)~\cite{FL} for optimizing monitoring policies, few have fully leveraged the programmability offered by P4 to jointly optimize the trade-off between control overhead and monitoring precision from both the data and control planes \cite{SONI2022103419}. 
% As far as is known, this is the first paper to propose an SDN monitoring framework that explicitly integrates programmable data plane mechanisms enabled by P4 with machine learning-based adaptive policy inference on the control plane. 
\rev{This paper presents the first SDN monitoring framework that integrates P4-enabled programmable data plane mechanisms with machine learning-based adaptive policy inference on the control plane.} 
This cross-plane coordination enables closed-loop, fine-grained monitoring with improved adaptability and efficiency. 
Based on this architectural principle, a framework named MAT4PM (Machine Learning-Guided Adaptive Threshold Control for P4-based Monitoring) is developed. 
It integrates event-driven monitoring logic implemented in the P4 data plane with predictive threshold adjustment conducted by control-plane machine learning models~\cite{MAT4PM}. 
The primary contributions are summarized as follows:

\begin{enumerate}
    
    \item A collaborative monitoring mechanism is designed and implemented by combining P4-based active reporting at the data plane with machine learning-driven policy tuning at the control plane. 
    Building upon the advantages of proactive schemes such as PPM (P4-based Proactive Monitoring) \cite{10623873}, the proposed approach eliminates the reliance on manually configured thresholds by dynamically predicting monitoring intervals through learned models with real-time traffic features as input. 
    The system establishes a closed-loop architecture characterized by data plane event triggering and control plane intelligent decision-making, thereby achieving high accuracy, low overhead, and strong adaptability under dynamic traffic scenarios.
    
    \item A data construction methodology is formulated to support optimal threshold learning and facilitate the generation of a high-quality labeled dataset. 
    Optimal thresholds are derived from offline simulations that minimize a jointly defined cost function, serving as ground truth labels for supervised learning. 
    Packet trace files (pcap)~\cite{pcap} captured from the BMv2 switch under various monitoring intervals are replayed to quantify estimation errors and communication overheads.
    
    \item A full prototype system is developed, comprising P4-based data plane logic, controller-side decision modules, and three mainstream machine learning models---eXtreme Gradient Boosting (XGBoost)~\cite{XGBoost}, Light Gradient Boosting Machine (LightGBM)~\cite{LightGBM}, and Support Vector Machine (SVM)~\cite{SVM}. 
    Experimental results demonstrate that, compared to fixed-threshold schemes, the proposed dynamic adjustment mechanism reduces the average estimation error from 9.1\% to 7.0\% and lowers the overall monitoring cost by 5.6\%. 
    In addition, the models exhibit low prediction latency (all below 5~ms) and minimal resource consumption (memory usage: 13.89 KB-36.24 KB), indicating the feasibility of real-time deployment in operational environments.
    
\end{enumerate}

% \rev{The terminology and abbreviations used throughout this paper are summarized in Table~\ref{tab:acronyms} for reference.} 
The remainder of this paper is organized as follows. Section~\ref{sec2} reviews background and related work, including fundamental concepts of programmable data planes and existing P4-based monitoring approaches. Section~\ref{sec3} details the architecture and core components of the proposed MAT4PM system. Section~\ref{sec4} presents experimental evaluations to validate the effectiveness of our method. Section~\ref{sec5} concludes the paper and discusses potential future research directions.

% \begin{table}[!t]
%     \centering
%     \caption{\rev{List of Acronyms and Definitions}}
%     \label{tab:acronyms}
%     \begin{tabular}{p{0.09\textwidth} p{0.35\textwidth}}
%         \toprule
%         \multirow{1}{=}{\textbf{Abbreviation}} & \textbf{Definition} \\
%         \midrule
%         \rev{API} & \rev{Application Programming Interface} \\
%         \rev{ASIC} & \rev{Application-Specific Integrated Circuit} \\
%         \rev{BMv2} & \rev{Behavioral Model version 2} \\
%         \rev{BPS} & \rev{Bits Per Second} \\
%         \rev{CI} & \rev{Confidence Interval} \\
%         \rev{CPU} & \rev{Central Processing Unit} \\
%         \rev{DLINT} & \rev{Deterministic Lightweight In-Band Network Telemetry} \\
%         \rev{DPDK} & \rev{Data Plane Development Kit} \\
%         \rev{FL} & \rev{Federated Learning} \\
%         \rev{FS-INT} & \rev{Flexible Sampling-based In-band Network Telemetry} \\
%         \rev{INT} & \rev{In-band Network Telemetry} \\
%         \rev{INTO} & \rev{In-band Network Telemetry Orchestration} \\
%         \rev{LightGBM} & \rev{Light Gradient Boosting Machine} \\
%         \rev{LB} & \rev{Load Balancing} \\
%         \rev{MAC} & \rev{Media Access Control} \\
%         \rev{MSE} & \rev{Mean Squared Error} \\
%         \rev{ML} & \rev{Machine Learning} \\
%         \rev{OVS} & \rev{Open vSwitch} \\
%         \rev{P4} & \rev{Programming Protocol-independent Packet Processor} \\
%         \rev{PCAP} & \rev{Packet Capture} \\
%         \rev{PDP} & \rev{Programmable Data Plane} \\
%         \rev{PLINT} & \rev{Probabilistic Lightweight In-Band Network Telemetry} \\
%         \rev{PPM} & \rev{P4-based Proactive Monitoring} \\
%         \rev{PPS} & \rev{Packets Per Second} \\
%         \rev{RAM} & \rev{Random Access Memory} \\
%         \rev{RL} & \rev{Reinforcement Learning} \\
%         \rev{RPC} & \rev{Remote Procedure Call} \\
%         \rev{Std. Dev.} & \rev{Standard Deviation} \\
%         \rev{SDN} & \rev{Software-defined Networking} \\
%         \rev{SL} & \rev{Supervised Learning} \\
%         \rev{SVM} & \rev{Support Vector Machine} \\
%         \rev{TCP} & \rev{Transmission Control Protocol} \\
%         \rev{TBSW} & \rev{Time-based Sliding Window} \\
%         \rev{UDP} & \rev{User Datagram Protocol} \\
%         \rev{VM} & \rev{Virtual Machine} \\
%         \rev{XGBoost} & \rev{eXtreme Gradient Boosting} \\
%         \bottomrule
%     \end{tabular}
% \end{table}
